\documentclass[../main.tex]{subfiles}
\graphicspath{{\subfix{../images/}}}
\begin{document}

\subsection{Experimental Setup}

We compared both approaches on synthetic data (50×50 surfaces, 10,000 points, various patterns) and real PLY data (floor.ply, 61,914 points). Configuration: 8×8 grid, subspace dimension 3, eight-neighbor computation, 50\% sliding window overlap.

\subsection{Results}

\subsubsection{Synthetic Data Performance}

On synthetic data, both approaches demonstrated effective pattern detection with comparable performance characteristics:
\begin{itemize}
    \item \textbf{Fixed Patches:} Max magnitude = 0.652, Mean = 0.494
    \item \textbf{Sliding Window:} Max magnitude = 0.770, Mean = 0.445
    \item \textbf{Resolution increase:} 3.6× (361 vs 100 analysis points)
\end{itemize}

The sliding window approach showed 18\% improvement in peak detection while maintaining similar mean values, confirming the theoretical advantage of overlapping analysis.

\subsubsection{Real-World Data Performance}

Real PLY file analysis revealed contrasting performance characteristics:
\begin{itemize}
    \item \textbf{Fixed Patches:} Max magnitude = 1.079, Mean = 0.977, Std = 0.074
    \item \textbf{Sliding Window:} Max magnitude = 0.888, Mean = 0.547, Std = 0.150
    \item \textbf{Resolution increase:} 3.5× (225 vs 64 analysis points)
\end{itemize}

\begin{figure}[tb]
  \centering
  \includegraphics[width=0.8\textwidth]{comparison_analysis.png}
  \caption{Comparative analysis of fixed patches vs sliding window on real PLY data. The sliding window approach demonstrates superior visual contrast and clearer anomaly localization despite lower absolute magnitude values.}
  \label{fig:comparison}
\end{figure}

\subsection{Key Findings}

\subsubsection{Data-Dependent Performance}

The most significant finding is that method effectiveness depends critically on data characteristics:

\textbf{Clean Synthetic Data:} Both methods perform well, with sliding window providing modest improvements in peak detection.

\textbf{Noisy Real Data:} Sliding window demonstrates superior practical performance through:
\begin{itemize}
    \item Better signal-to-background ratio (wider dynamic range: 0.22-0.89 vs 0.73-1.08)
    \item Clearer visual contrast for anomaly identification
    \item Reduced false positive rates through noise averaging
\end{itemize}

\subsubsection{Visual Contrast Analysis}

Despite lower absolute magnitude values on real data, the sliding window approach provides enhanced background suppression with preserved peak signals, creating clearer visual contrast for anomaly identification. The overlapping analysis creates smoother gradients and enables more precise anomaly localization through higher spatial resolution.


\end{document}